\documentclass[10pt,a4paper]{article}
\usepackage[T1]{fontenc}
\usepackage[margin=22mm]{geometry}
\usepackage{lmodern,amsmath,amssymb,amsthm,mathtools,microtype}
\usepackage[hidelinks]{hyperref}
\newtheorem{theorem}{Theorem}
\title{Factorial one-sided count certificates\newline on the odd reciprocal ES spectrum}
\author{The Clankers}
\date{14 September 2026 --- research preprint}
\begin{document}\maketitle
\begin{abstract}
The exact denominator defects of the Erd\H{o}s--Straus divisor atlas are odd.
Explicit rational interpolation polynomials therefore separate successes from
failures with factorial uniform error and only linear growth of the absolute
coefficient sum. Preserving the original state counts gives simultaneous
one-sided certificates for every canonical row defined here, for primes
$p\equiv1\pmod4$, at moment degree asymptotic to
$\log p/\log\log p$. The finite coefficient-error and split-compression maps
are explicit. The theorem controls recovery of the count; it does not prove
that every such prime's count is positive. No historical-priority claim is made.
\end{abstract}

\paragraph{Exact arithmetic domain.}
For prime $p=4m+1$, retain every $m+1\le a\le2m$, $R=4a-p$,
$a=\prod_iq_i^{e_i}$, $-e_i\le\beta_i\le e_i$, $u=\prod_iq_i^{e_i+\beta_i}$,
and the two separate channel tags $E,M$. Their gates and defects are
\[
 g_E=4u+1,\qquad g_M=u+a,\qquad D=R/\gcd(R,g_c).
\]
The original positive rational triples are
\[
 E:\ \left(a,\frac{pa+a^2/u}{R},\frac{pa+p^2u}{R}\right),\qquad
 M:\ \left(a,\frac{p(a+a^2/u)}R,\frac{p(a+u)}R\right).
\]
Multiplication gives $(RY-pa)(RZ-pa)=p^2a^2$, proving their reciprocal sum is
$4/p$. Since $\gcd(pa,R)=1$ and $4a\equiv p\pmod R$, both nonfirst reduced
denominators are exactly $D$. Thus the hit indicator is $\mathbf1_{D=1}$,
and every failure has odd $D\ge3$. The inherited complete first-half atlas
retains the Type I/II coordinate conventions \cite{ET,Prior}. The original
inverse is $u=a^2/(RY-pa)$ for E and $u=pa^2/(RY-pa)$ for M, with its tag.

For the countable disjoint union of these labelled rows use counting
$\ell^2$, $Xe_\alpha=D_\alpha^{-1}e_\alpha$, and
$\Pi e_\alpha=\mathbf1_{D_\alpha=1}e_\alpha$. The spectral accumulation
point zero is not an absent state. The exact ambient compact set is
\[
 \Sigma=\{0,1,1/3,1/5,1/7,\ldots\}.
\]

\begin{theorem}[Discrete signed interpolation]
Define the degree-$(j+1)$ polynomial
\[
 \Phi_j(x)=x\prod_{\nu=1}^j\frac{(2\nu+1)x-1}{2\nu},\qquad j\ge0.
 \tag{1}
\]
Its value is one at 1. At $D=2k+1>1$ its value is zero for $1\le k\le j$;
otherwise it is
\[
 \Phi_j(1/D)=(-1)^j\frac{\binom{k-1}{j}}{(2k+1)^{j+1}}.
 \tag{2}
\]
Consequently on $\Sigma$,
\[
 \Phi_{2r+1}\le\mathbf1_{\{1\}}\le\Phi_{2r},\qquad
 \|\Phi_j-\mathbf1_{\{1\}}\|_\Sigma
 \le\frac1{2^j j!(2j+3)}.
 \tag{3}
\]
The even sequence decreases and the odd sequence increases to the indicator.
The coefficient $\ell^1$ norm of $\Phi_j$ is exactly $j+1$.
\end{theorem}
\begin{proof}
Each factor of (1) is $(\nu-k)/(\nu D)$ at $1/D$, giving (2).
For $k>j$, its absolute value is at most $1/(2^j j!D)$, and $D\ge2j+3$.
The two additional factors in the same-parity comparison have product in
$[0,1]$ on each remaining failed node; newly reached nodes become zero.
This proves the signs and monotonicity. All nonzero coefficients alternate;
their absolute sum is $|\Phi_j(-1)|=\prod_{\nu=1}^j(\nu+1)/\nu=j+1$.
\end{proof}

\newpage
\paragraph{Interpolation and exact error.}
The full evaluation map on $0,1,1/3,\ldots,1/(2j+1)$ is a rational linear
isomorphism. Its inverse is the Lagrange formula
\[
 (v_i)_i\longmapsto\sum_i v_i\prod_{h\ne i}\frac{x-z_h}{z_i-z_h}.
\]
Thus (1) is one explicitly selected inverse column of a classical map
\cite[\S3.3(i), equations 3.3.1--3.3.3]{DLMF}. Its values outside these nodes
are not asserted to have the signs in (3); the discrete spectrum is essential.

There is a sharp finite alternative to the bound in (3). For $j\ge1$ the
maximum in (2) occurs among the integers
\[
 \max\{j+1,\lfloor1+j(j+2)/2\rfloor\}\le k\le
 \max\{j+1,\lceil j+j(j+2)/2\rceil\}.
 \tag{4}
\]
Indeed the logarithmic derivative of the positive expression in (2) has the
sign of $\sum_{\nu=1}^j(2\nu+1)/(k-\nu)-2$, a strictly decreasing function.
Its crossing lies between $1+j(j+2)/2$ and $j+j(j+2)/2$ by bounding every
denominator between $k-j$ and $k-1$. This proves (4), including the endpoint
restriction. Exact evaluation gives errors $1/3,2/49,10/2197,4/9261$ at
$j=0,1,2,3$, attained respectively at $D=3,7,13,21$.

\paragraph{Original mass, not a probability replacement.}
Let $W_p$ be the size of a full canonical row and $H_p$ its success count.
An elementary bound is
\[
 W_p\le B_p=\min\{6m^2,4m(1+\lfloor\log_2(2m)\rfloor)^2\}.
 \tag{5}
\]
For the first expression use $\tau(a^2)\le2a-1$. For the second use the
prime-power inequality $2e+1\le(e+1)(e+2)/2$, hence $\tau(a^2)\le d_3(a)$.
Then $\sum_{a\le N}d_3(a)\le N(\sum_{u\le N}1/u)^2$, and dyadic grouping
bounds the harmonic sum by $1+\lfloor\log_2N\rfloor$. Double for both channels.

\begin{theorem}[Simultaneous original-count certificates]
For accuracy index $n\ge0$, choose the least even $j=j(p,n)$ with
\[
 2^j j!(2j+3)\ge8\cdot3^nB_p.
\]
Set $Q_j(p)=\sum_{\alpha\in\mathcal A_p}\Phi_j(D_\alpha^{-1})$. Then
\[
 \boxed{H_p-\frac1{16\cdot3^n}\le Q_{j+1}(p)\le H_p
 \le Q_j(p)\le H_p+\frac1{8\cdot3^n}.}
 \tag{6}
\]
In particular $\lceil Q_{j+1}(p)\rceil=\lfloor Q_j(p)\rfloor=H_p$, and
$H_p>0$ if and only if $Q_{j+1}(p)>0$. The largest required moment index is
$j+2=(1+o(1))\log p/\log\log p$ for fixed $n$.
\end{theorem}
\begin{proof}
Sum (3) with the original unit masses and apply (5).
The ratio between consecutive error denominators is
$2(j+1)(2j+5)/(2j+3)>2$. This proves (6) and its rounding consequence.
The estimate $\log(2^j j!(2j+3))=j\log j+O(j)$ follows by integrating
$\log x$ above and below its integer sum. Minimality among even integers
changes it by $O(\log j)$; $\log B_p=\log p+O(\log\log p)$ proves the degree.
\end{proof}
The inequality $Q_{j+1}(p)>0$ is not established for every prime by this
convergence theorem. It is the now-explicit arithmetic lower bound to prove.
On the original prime counting space, the number operator $N=\operatorname{diag}H_p$
retains its maximal domain $\sum H_p^2|v_p|^2<\infty$.
The two operators in (6) have that same domain and bounded differences from
$N$ with the displayed uniform constants.

\newpage
\paragraph{Certified moment inputs.}
Write $M_k(p)=\sum_\alpha D_\alpha^{-k}$. Each $Q_j$ is an explicit rational
linear combination of these original moments. If every moment through degree
$j+2$ has certified absolute error at most
\[
 \delta=\frac1{16(j+2)3^n},
\]
the total evaluation error is at most $1/(16\cdot3^n)$ for either filter.
Form the lower endpoint by subtracting this allowance and the upper by adding
it to the approximate sums. Their width is at most $7/(16\cdot3^n)<1$ and they
contain the integer $H_p$. This is an absolute error statement. It does not
claim that computing the original moments to this accuracy has low complexity.

For any finite original selection, the particularly simple quartic gives
\[
 \Phi_3(x)=\frac{105x^4-71x^3+15x^2-x}{48}\le\mathbf1_{\{1\}}(x).
\]
At one shell put $S_k=\sum_\alpha\gcd(R,g_\alpha)^k$. The unconditional
sufficient test is therefore
\[
 \boxed{105S_4-71RS_3+15R^2S_2-R^3S_1>0\quad\Longrightarrow\quad H>0.}
 \tag{7}
\]
A separate accompanying note supplies a hard-prime shell at which (7) succeeds
while the optimal aggregate first-degree signed moment test fails. Neither (7)
nor first-degree adequacy is asserted necessary for arbitrary positive shells.

\paragraph{Split-fibre and operator return.}
For $e=f+b$ the addition map $(\xi,\eta,c)\mapsto(\xi+\eta,c)$ has complete
fibre
\[
 \max(-f_i,\beta_i-b_i)\le\xi_i\le\min(f_i,\beta_i+b_i),
 \quad\eta=\beta-\xi.
\]
Its size $n(\beta)$ is the product of those interval lengths. Assigning each
occurrence weight $1/n(\beta)$ preserves every $M_k$ and $Q_j$ exactly by
finite fibre summation. All channel labels and original exponent vectors
remain; a residue-support subgroup does not replace the exponent box.

In the completed occurrence space with those weights, value duplication $U$
has adjoint equal to fibre averaging, and $U^*U=I$. For $A=UXU^*$ on the full
occurrence space, every polynomial satisfies
\[
 f(A)=Uf(X)U^*+f(0)(I-UU^*).
\]
Decompose into $\operatorname{im}U$ and its orthogonal complement to prove it.
Because $\Phi_j(0)=0$,
\[
 \boxed{\Phi_j(UXU^*)=U\Phi_j(X)U^*.}
 \tag{8}
\]
The return on the image uses $U^*$; the kernel directions stay marked.
The raw repeated diagonal trace is different and multiplies original states
by their fibre sizes. For the product in (1) without the leading $x$, its
constant term would instead contribute
$(-1)^j(2^jj!)^{-1}(I-UU^*)$.

\paragraph{Verification and scope.}
The accompanying standard-library verifier uses rational arithmetic and checks
all signs, coefficient sums, interpolation inverse columns, sharp finite
windows, original split fibres and the stated operator comparisons. It checks
the row schedule through $p\le1500$ and on explicit growing-mass fixtures.
The infinite theorems are proved above, not inferred from that finite replay.
Original normalized ES coordinates recover a witness after a positive sum
identifies a successful finite source label. Moments alone do not identify
which state supplied it. Supported coefficient zero is not external absence.

\begin{thebibliography}{3}\small
\bibitem{DLMF} National Institute of Standards and Technology. \emph{Digital Library
of Mathematical Functions}, \S3.3(i), equations 3.3.1--3.3.3.
\bibitem{ET} Elsholtz, C., \& Tao, T. (2013). Counting the number of solutions to
the Erd\H{o}s--Straus equation on unit fractions. \emph{J. Aust. Math. Soc., 94}(1),
50--105, \S2, Propositions 2.2 and 2.6.
\bibitem{Prior} The Clankers (2026). \emph{Uniform infinite-array control for marked
Erd\H{o}s--Straus selectors}, 14 September, Lemma 1 and Theorem 2. Complete
source/hash identity and the preceding SplitZero integration are in the companion ledger.
\end{thebibliography}
\end{document}
